% !TeX root = ../memoire.tex

\chapter*{Introduction}
\addcontentsline{toc}{chapter}{Introduction}
\markboth{Introduction}{Introduction}

\noindent\hfill
\begin{minipage}{0.72\textwidth}
	\small
	\itshape
	\enquote{L’histoire de la Faculté de Médecine est désormais possible : il est éminemment souhaitable qu’elle soit entreprise. […] Combien de thèses ou de mémoires pourraient être tirés des archives de la Faculté et des dépôts qui les complètent ? […] Nous souhaiterions, notamment, qu’une série de travailleurs se rencontrât pour étudier les maîtres de l’ancienne Faculté, tant au point de vue de leur vie qu’au point de vue de leurs œuvres. Les archives fourniraient une matière inépuisable de renseignements inédits ; la Bibliothèque donnerait, en guise de complément, les éléments imprimés. Quel livre intéressant ne ferait-on pas, par exemple, sur Laurent Joubert dont nos documents évoquent la figure si vivante et dont les traités copieux tiennent une si large place ? […] Mais, qu’il s’agisse des médecins, des chirurgiens ou des apothicaires, les sujets de tout ordre surgissent en foule, et il suffit de feuilleter l’Inventaire […] pour apprécier les ressources que cet ensemble de documents inexplorés réserve à la bonne volonté des travailleurs à venir.}
	
	\begin{flushright}
		\normalfont\small
		Joseph Calmette, \emph{Cartulaire de l’Université de Montpellier}, t.~II, 1912, p.~V--VII.
	\end{flushright}
\end{minipage}

\vspace{1.5em}

En 1903, Joseph Calmette, archiviste paléographe, est chargé par la Faculté de médecine de Montpellier de reprendre le classement de ses archives anciennes et d'en établir l'inventaire. Mené jusqu'en 1908 sur quelque vingt-cinq mille documents, ce travail conduit notamment à la redécouverte du premier inventaire général connu des archives de la Faculté, dressé en 1583 à la suite d'une entreprise de rassemblement de titres alors dispersés\footcite[p.~I et XII, vues~20 et 31/1131]{CalmetteCartulaire1912}. Les résultats du classement sont publiés en 1912 dans le second tome du \emph{Cartulaire de l'Université de Montpellier}, vingt-deux ans après un premier volume consacré à l'édition des principaux actes relatifs à l'Université de médecine depuis le XII\textsuperscript{e} siècle. L'inventaire établi par Calmette, qui demeure aujourd'hui l'instrument de référence pour le fonds ancien, y est publié aux côtés de celui de 1583, réunissant dans un même volume deux états des archives séparés par plus de trois siècles.

Ces archives sont celles d'une institution dont l'histoire remonte à la constitution de l'enseignement médical montpelliérain en \emph{universitas medicorum} par les statuts de 1220, et que l'Université présente aujourd'hui comme \enquote{la plus ancienne école universitaire de médecine au monde toujours en exercice}\footcite{UM800Ans2020}. Au fil de ses transformations, les documents ont circulé entre officiers et professeurs, été déplacés, dispersés puis réunis et reclassés ; l'inventaire de 1583 mentionne déjà les titres que le chancelier et les docteurs ont pu \enquote{estre recouvert des egarez}\footcite[p.~XII, vue~31/1131]{CalmetteCartulaire1912}. Les productions administratives des XIX\textsuperscript{e} et XX\textsuperscript{e} siècles, des fonds privés et des collections iconographiques ont par la suite rejoint les ensembles plus anciens et fait l'objet de traitements successifs. Les Archives historiques conservent ainsi, dans leur composition comme dans leur organisation, la trace des institutions qui les ont produites ou recueillies et des opérations qui ont accompagné leur transmission.

Cette relation étroite entre les archives et l'institution qui les a produites prend aujourd'hui une importance particulière. Les Archives historiques sont toujours conservées dans le bâtiment historique de la Faculté de médecine et leur traitement est assuré au sein même de l'Université de Montpellier. Relevant du régime des archives publiques, elles ont pourtant en principe vocation à être versées au service public d'archives compétent ; leur maintien à la Faculté suppose par conséquent une dérogation à cette obligation de versement\footcite[art.~L.~211-4, L.~212-4, I, et R.~212-12]{CodePatrimoine}. Depuis 2023, l'Université prépare cette demande avec les Archives départementales de l'Hérault dans le cadre d'une campagne associant classement, reprise et création d'instruments de recherche, conservation préventive, restauration, numérisation et valorisation. L'enjeu dépasse ainsi la seule localisation matérielle des documents : il s'agit d'organiser durablement, au sein d'un établissement d'enseignement supérieur, les conditions dans lesquelles des archives publiques peuvent être conservées, décrites, communiquées et valorisées tout en demeurant auprès de l'institution dont elles documentent l'histoire.

La diffusion numérique prend place dans ce contexte. Les Archives historiques disposent déjà d'instruments de recherche élaborés à différentes périodes et de ressources diffusées dans plusieurs environnements ; la campagne en cours en crée de nouveaux et reprend progressivement certaines descriptions existantes. Leur mise en ligne soulève dès lors une question qui ne se réduit pas au choix d'un site ou d'un logiciel : elle engage plusieurs fonctions --- médiation, signalement archivistique, diffusion des reproductions --- et plusieurs services de l'Université, auxquels s'ajoutent les institutions et réseaux documentaires extérieurs susceptibles de relayer ou d'enrichir cet accès. La conservation \emph{in situ} conduit ainsi à interroger la place des Archives historiques dans un environnement plus large que celui de la seule Faculté, à l'échelle de l'Université puis de l'enseignement supérieur et de la recherche.

\paragraph{Cadre de l'étude}

Le stage à l'origine de ce mémoire a été réalisé au sein de la Mission Archives historiques de l'UFR de médecine de Montpellier, sous la responsabilité de Sophie Dikoff, cheffe de la Mission et tutrice du stage, avec Emma Bertrand, assistante archiviste. La commande formulée en avril 2026 portait sur l'étude et la préfiguration d'un portail consacré aux Archives historiques, destiné à donner accès aux fonds, à leurs instruments de recherche et à des contenus de valorisation. Les premiers échanges avec la Direction des systèmes d'information et du numérique ont cependant montré qu'aucune solution immédiatement mobilisable ne permettait d'envisager sa mise en production à court terme. Le stage a dès lors pris la forme d'une étude préalable des conditions documentaires, techniques et institutionnelles de cette diffusion, accompagnée de prototypes destinés à éprouver concrètement certaines possibilités.

L'enquête a progressivement dépassé le périmètre du seul portail pour examiner les différents dispositifs susceptibles d'intervenir dans l'accès aux fonds, les services qui en ont la responsabilité et les circulations possibles entre descriptions archivistiques, reproductions numériques et contenus éditoriaux. L'étude de solutions employées dans d'autres établissements et les échanges menés avec leurs responsables ont complété l'analyse de l'environnement montpelliérain. Ce travail doit aboutir à une feuille de route qui n'a pas vocation à retenir par avance une architecture définitive, mais à délimiter les différentes trajectoires qui pourront être suivies par la Mission en fonction des choix techniques et institutionnels ultérieurs, en précisant pour chacune les étapes, les prérequis et les acteurs concernés.

\paragraph{Sources et méthode}

L'enquête historique s'est appuyée sur les archives conservées à la Faculté et sur leurs instruments de recherche, complétés par les documents identifiés aux Archives départementales de l'Hérault, aux Archives municipales de Montpellier et dans d'autres institutions de conservation. Les statuts, délibérations, récépissés, anciens inventaires, plans de classement et dossiers administratifs ont notamment été confrontés aux indications données dans les éditions de sources, en particulier dans le second tome du \emph{Cartulaire}\footcite{CalmetteCartulaire1912}. Leur croisement permet de restituer la constitution et les réorganisations successives des archives tout en tenant compte des lacunes et des incertitudes de leur transmission.

Cette enquête historique rejoint directement les problématiques archivistiques qui structurent le mémoire. Les relations entre histoire et archives ont en effet donné lieu, depuis plusieurs décennies, à des réflexions qui ont interrogé parallèlement les effets des processus d'archivage sur les fonds et la manière dont les historiens appréhendent les archives comme objets de recherche. Ces approches, longtemps développées dans des cadres en partie distincts, invitent à rapprocher davantage les questionnements des historiens et ceux des archivistes\footcite{Poncet2019ArchivesHistoire}. Cette articulation prend ici une dimension concrète : l'histoire des institutions et de leurs archives permet d'éclairer l'organisation des fonds, tandis que leur traitement archivistique conditionne les formes sous lesquelles ils peuvent aujourd'hui être connus et mobilisés.

L'étude des conditions de diffusion a combiné l'examen des instruments existants, des essais techniques et des échanges avec les acteurs concernés. Plusieurs instruments de recherche ont été comparés, encodés ou repris en EAD puis testés dans SimplEAD afin d'examiner les possibilités de publication et de signalement. La dimension éditoriale du futur site a été éprouvée par un premier prototype réalisé localement en HTML, CSS et JavaScript, puis par un second dans le bac à sable WordPress de l'Université.

Les essais consacrés aux données archivistiques ont également porté sur l'instance de Flora mise à disposition avant sa mise en production : ses possibilités de description et d'export ont été examinées à partir de données des Archives historiques, tandis que des scripts ont été développés pour analyser et corriger les anomalies rencontrées dans l'export EAD. Les possibilités offertes par Calames et FranceArchives pour le signalement, ainsi que par IIIF et Mirador pour la diffusion et la consultation des reproductions, ont été étudiées à partir de leur documentation. Le fonctionnement de Foli@, qui articule au sein de l'Université Flora, Drupal et la diffusion des images, a fait l'objet d'échanges réguliers avec le Service de coopération documentaire interuniversitaire, permettant d'examiner les conditions dans lesquelles les Archives historiques pourraient s'inscrire dans une infrastructure documentaire déjà existante.

Ces expérimentations ont été complétées par des entretiens et des retours d'expérience auprès des services de l'Université, des Archives départementales et municipales de Montpellier et d'établissements confrontés à des problématiques comparables. Le périmètre demeure celui d'une étude préalable : la campagne de traitement se poursuit et plusieurs paramètres techniques ou institutionnels restent susceptibles d'évoluer. Les prototypes, tests et scénarios permettent précisément de documenter ces choix avant qu'ils ne soient arrêtés et de borner les différentes voies qui pourront être suivies par la Mission.

À partir de ce terrain, le mémoire interroge les formes que peut prendre l'accès numérique à des archives maintenues auprès de l'institution qui les a produites : \emph{comment organiser la diffusion numérique d'archives historiques conservées au sein de l'université qui les a produites, tout en conciliant leur ancrage institutionnel, leur intelligibilité et leur inscription dans les réseaux documentaires et scientifiques ?}

L'étude revient d'abord sur les institutions qui ont successivement produit, conservé et organisé ces archives, afin de comprendre la constitution des fonds, les classements dont ils ont fait l'objet et les conditions juridiques dans lesquelles se poursuit aujourd'hui leur conservation \emph{in situ}. Elle distingue ensuite les différentes fonctions que recouvre leur diffusion numérique --- médiation éditoriale, publication et \gls{signalement} des instruments de recherche, diffusion des reproductions --- pour examiner les solutions susceptibles d'y répondre, leurs limites et leurs possibilités d'articulation ; cette analyse permet de circonscrire les différentes voies qui resteront ouvertes à la Mission en fonction des arbitrages techniques et institutionnels ultérieurs. La dernière partie change enfin d'échelle pour replacer les Archives historiques dans l'écosystème patrimonial et scientifique de l'\gls{esr}. Elle examine d'abord les relations qui peuvent être établies avec les collections bibliographiques, muséales et scientifiques de l'Université sans uniformiser leurs descriptions, puis l'inscription des fonds dans des réseaux et des ensembles documentaires interinstitutionnels. Elle envisage enfin la constitution, à partir de certaines sources, de \glspl{corpus-numerique} et de données de recherche, en interrogeant les transformations successives des documents et le maintien de leur \gls{provenance-donnees}.